\chapter{Personalización de tu Proyecto}

Este capítulo explica la forma recomendada de ampliar el firmware en
\texttt{MICSBOL/ESP32\_BT\_Controller}. El repositorio está organizado de manera que:

\begin{itemize}
    \item \textbf{El mapeo de pines está centralizado} en \texttt{pins.h} y depende del modo de compilación seleccionado
    mediante \texttt{PROJECT\_MODE}.
    \item \textbf{Las lecturas de hardware} (ADC/GPIO/I\textsuperscript{2}C) se recolectan en módulos dedicados como
    \texttt{input\_hw.*} (lecturas ADC/digitales) y \texttt{i2c\_sensors.*}.
    \item \textbf{Los paquetes de telemetría} se construyen y se transmiten en \texttt{telemetry.cpp}, pero los \emph{valores}
    provienen de un módulo proveedor separado, \texttt{telemetry\_source.*}.
    \item \textbf{La expansión opcional de GPIO} (MCP23017) se gestiona mediante \texttt{mcp\_io.*} y el mapeo dependiente
    del modo en \texttt{pins.h}.
\end{itemize}

\section{Mapeo de Pines (pins.h) y Modos de Proyecto}

El repositorio define dos modos de cableado en \texttt{pins.h}:

\begin{itemize}
    \item \textbf{Modo directo ESP32} (\texttt{PROJECT\_MODE == FULL\_RC\_MODE}): la mayoría de E/S usa pines GPIO del ESP32.
    \item \textbf{Modo expansión MCP} (\texttt{PROJECT\_MODE == FULL\_RC\_MODE\_MCP}): la E/S digital se amplía usando un
    MCP23017 por I\textsuperscript{2}C, mientras que las salidas PWM se mantienen en el ESP32.
\end{itemize}

\subsection{Bus I\textsuperscript{2}C (reservado / bus MCP)}
En ambos modos, \texttt{pins.h} reserva el bus I\textsuperscript{2}C del ESP32 en:

\begin{itemize}
    \item \texttt{PIN\_I2C\_SDA = 21}
    \item \texttt{PIN\_I2C\_SCL = 22}
\end{itemize}

En modo MCP, estos pines se usan para la comunicación ESP32 $\leftrightarrow$ MCP23017 (y también pueden compartirse con otros
sensores I\textsuperscript{2}C si el bus está diseñado correctamente).

\subsection{Salidas PWM (6 en total)}
En ambos modos, \texttt{pins.h} define seis salidas PWM en el ESP32:

\begin{itemize}
    \item \texttt{PIN\_PWM\_STICK\_LX = 18}
    \item \texttt{PIN\_PWM\_STICK\_LY = 19}
    \item \texttt{PIN\_PWM\_STICK\_RX = 23}
    \item \texttt{PIN\_PWM\_STICK\_RY = 25}
    \item \texttt{PIN\_PWM\_KNOB\_L = 26}
    \item \texttt{PIN\_PWM\_KNOB\_R = 27}
\end{itemize}

Estos canales suelen usarse para servos, entradas PWM de drivers de motor u otras salidas de tipo analógico.

\subsection{Salidas de dirección de motor (2)}
Ambos modos definen dos salidas de dirección pensadas para drivers de motor:

\begin{itemize}
    \item \texttt{PIN\_MD\_LEFT = 12} \quad (nota en el repositorio: \emph{mantener en LOW al arranque})
    \item \texttt{PIN\_MD\_RIGHT = 15} \quad (nota en el repositorio: \emph{mantener en LOW al arranque})
\end{itemize}

\paragraph{Nota de seguridad al arranque}
Los comentarios en \texttt{pins.h} advierten que estos pines deben mantenerse en LOW durante el arranque. Al personalizar,
evita resistencias pull-up externas o circuitos que lleven estos pines a HIGH durante el reset.

\subsection{Entradas analógicas (ADC1 seguras)}
El repositorio usa explícitamente pines de ADC1 para entradas analógicas (recomendado cuando se usa Bluetooth).
Sin embargo, los números de pin cambian según el modo:

\paragraph{Modo directo ESP32 (\texttt{FULL\_RC\_MODE})}
\begin{itemize}
    \item \texttt{PIN\_NUMERIC1 = 34}
    \item \texttt{PIN\_NUMERIC2 = 35}
    \item \texttt{PIN\_ANALOG\_IND1 = 36}
    \item \texttt{PIN\_ANALOG\_IND2 = 39}
\end{itemize}

\paragraph{Modo expansión MCP (\texttt{FULL\_RC\_MODE\_MCP})}
\begin{itemize}
    \item \texttt{PIN\_NUMERIC1 = 32}
    \item \texttt{PIN\_NUMERIC2 = 33}
    \item \texttt{PIN\_ANALOG\_IND1 = 34}
    \item \texttt{PIN\_ANALOG\_IND2 = 35}
\end{itemize}

\paragraph{Implicación práctica}
Si cambias de modo de proyecto, \textbf{debes} revisar de nuevo el cableado de los sensores analógicos, porque los pines ADC
para numéricos/indicadores cambian entre modos.

\subsection{Entradas digitales (indicadores) en modo directo}
En modo directo ESP32, \texttt{pins.h} define cuatro entradas digitales para indicadores:

\begin{itemize}
    \item \texttt{PIN\_IND1 = 4}
    \item \texttt{PIN\_IND2 = 5}
    \item \texttt{PIN\_IND3 = 2}
    \item \texttt{PIN\_IND4 = 0} \quad (nota del repositorio: \emph{cuidado: evitar pull-up al arranque})
\end{itemize}

\paragraph{Advertencia de arranque en GPIO0}
Como \texttt{PIN\_IND4} usa GPIO0 en modo directo, evita un cableado que fuerce GPIO0 a HIGH/LOW durante el reset de una forma que
ponga al ESP32 en un modo de arranque incorrecto.

\subsection{Salidas de interruptores (modo directo)}
En modo directo, el repositorio define salidas de interruptores en pines del ESP32:

\begin{itemize}
    \item \texttt{PIN\_SW1 = 16}
    \item \texttt{PIN\_SW2 = 17}
    \item \texttt{PIN\_SW3 = 13}
    \item \texttt{PIN\_SW4 = 14}
\end{itemize}

\paragraph{Importante}
En el estado actual de \texttt{pins.h}, solo se definen \texttt{PIN\_SW1} hasta \texttt{PIN\_SW4} en modo directo,
aunque otras partes del sistema se refieren a ``Switches 1--6'' de forma conceptual. Si tu proyecto necesita SW5/SW6 en modo directo,
debes asignar pines GPIO adicionales en \texttt{pins.h} y actualizar la capa de salidas en consecuencia.

\subsection{Salidas de pulso (botones de evento Android)}
Las salidas de pulso están pensadas para acciones momentáneas disparadas por botones de evento en Android.

\paragraph{Modo directo ESP32}
\begin{itemize}
    \item \texttt{PIN\_PULSE1 = 32}
    \item \texttt{PIN\_PULSE2 = 33}
\end{itemize}

\paragraph{Modo expansión MCP}
En modo MCP, \texttt{pins.h} asigna cuatro salidas de pulso en GPIO del ESP32:

\begin{itemize}
    \item \texttt{PIN\_PULSE1 = 4}
    \item \texttt{PIN\_PULSE2 = 5}
    \item \texttt{PIN\_PULSE3 = 13}
    \item \texttt{PIN\_PULSE4 = 14}
\end{itemize}

\section{Agregar Sensores (Fuentes de Telemetría)}

La forma más limpia de agregar sensores es alimentar sus mediciones en los canales de telemetría existentes. El API del proveedor
se declara en \texttt{telemetry\_source.h}:

\begin{itemize}
    \item \textbf{Paneles:} \verb|getPanelLeft()| y \verb|getPanelRight()| devuelven valores de 16 bits.
    \item \textbf{Indicadores:} \verb|getIndicatorAnalog()|, \verb|getIndicatorBattery()| y
    \verb|getIndicatorDigitalMask()| devuelven valores de 8 bits.
    \item \textbf{Gráficas (plots):} \verb|getPlot1()|, \verb|getPlot2()| y \verb|getPlot3()| devuelven valores de 8 bits.
\end{itemize}

\subsection{Sensores analógicos}
Flujo de trabajo típico:
\begin{enumerate}
    \item Elige un pin ADC desde la sección del modo correcto en \texttt{pins.h}.
    \item Implementa el escalado/conversión en \texttt{input\_hw.cpp} (consulta funciones como \verb|hwReadPlot1()|,
    \verb|hwReadIndicatorAnalog()|, etc., declaradas en \texttt{input\_hw.h}).
    \item Devuelve el valor procesado en \texttt{telemetry\_source.cpp} mediante \verb|getPlotX()| / \verb|getIndicator...()|.
\end{enumerate}

\subsection{Sensores I\textsuperscript{2}C}
El repositorio ofrece helpers I\textsuperscript{2}C en \texttt{i2c\_sensors.*} (por ejemplo, \verb|i2cReadTemp()| y
\verb|i2cReadIMU(...)| en \texttt{i2c\_sensors.h}). Para agregar otro dispositivo:
\begin{enumerate}
    \item Conecta SDA/SCL a \texttt{PIN\_I2C\_SDA}/\texttt{PIN\_I2C\_SCL}.
    \item Agrega una función de lectura en \texttt{i2c\_sensors.cpp/.h}.
    \item Publica el valor a través de \texttt{telemetry\_source.cpp}.
\end{enumerate}

\section{Uso de MCP23017 (Expansión GPIO)}

En modo MCP, \texttt{pins.h} documenta el mapeo digital previsto:

\begin{itemize}
    \item \textbf{GPA0--GPA7} $\rightarrow$ 8 indicadores digitales (ENTRADAS)
    \item \textbf{GPB0--GPB5} $\rightarrow$ 6 salidas de interruptores
    \item \textbf{GPB6--GPB7} $\rightarrow$ 2 salidas libres para que las uses
\end{itemize}

La interfaz del MCP en \texttt{mcp\_io.h} ofrece:

\begin{itemize}
    \item \verb|mcpInit()|
    \item \verb|mcpReadDigitalInputs()| (lee el puerto A)
    \item \verb|mcpWriteOutput(...)| y \verb|mcpWritePortB(...)| (escribe el puerto B)
\end{itemize}

\section{Lista Práctica al Personalizar}

\begin{itemize}
    \item \textbf{Comienza siempre revisando \texttt{pins.h} para el modo seleccionado}: los pines analógicos y de pulso cambian entre modos.
    \item \textbf{Ten cuidado con pines sensibles al arranque}: GPIO0 y GPIO12/15 tienen advertencias explícitas relacionadas con el arranque en el repositorio.
    \item \textbf{Si necesitas más E/S digital, prefiere el modo MCP}: amplía de forma limpia indicadores y salidas de interruptores mediante I\textsuperscript{2}C.
    \item \textbf{Cambia una cosa a la vez}: actualiza el cableado, luego actualiza \texttt{pins.h}/funciones de lectura, y después prueba telemetría y salidas.
    \item \textbf{Recomendación al flashear}: carga el firmware con el ESP32 \textbf{totalmente desconectado del hardware externo} (motores, servos, relés, sensores y especialmente cualquier cableado a pines de arranque/strapping). Esto evita fallos de carga por ruido eléctrico, caídas de tensión o por forzar niveles lógicos incorrectos en pines críticos durante el arranque.
    \item \textbf{Si falla la carga}: primero desconecta todo lo externo, prueba con otro cable USB/u otro puerto, y si tu placa lo requiere mantén presionado el botón \textbf{BOOT} al iniciar la carga. Reconecta el hardware externo \textbf{solo después} de confirmar que el firmware arranca correctamente.
\end{itemize}